\scanpage{30}
\mspage{24}

et les \deleted{\unsure{cond.}} inclusions stipulées dans $2^\circ$, b$'$)
assurent qu'on peut bel et bien trouver
$TC$ et $C$, quitte à les prendre les
plus grandes possibles $TC=W^{\mathrm{univ}}$, $C=\mathrm{Cof}_W$
--- auquel cas, d'ailleurs, on aura
\[
 TC=C\cap W^{\mathrm{univ}}_{\unsure{-}},
\]
i.e. on a une \inserted{pré-$Q$-}catégorie de \deleted{\unsure{Qu}}\unsure{modèles} à
g. Mais il s'agit d'assurer maintenant,
par contre, que
\[
 TF=F\cap W_{\mathrm{univ}}\unsure{)|3)|?},
\]
\deleted{Prenons}
ce qui implique, on vient de le voir \unsure{$*$} (en
notant $(W^{\mathrm{univ}})_*\subset F$) que
\[
 \begin{aligned}
 \overline{(W^{\mathrm{univ}})_*\cap W_{\mathrm{univ}}}
 &\subset\overline{F\cap W_{\mathrm{univ}}}=\overline{TF}={}&TF\\[-2pt]
 &&\cap\\[-2pt]&&W_{\mathrm{univ}}.
 \end{aligned}\tag*{($*$)}
\]

Inversement, si $(*)$ est satisfait, prenons
\[
 \begin{aligned}
 F&=(W^{\mathrm{univ}})_*\quad\text{le plus petit possible}\\
 &\text{i.e.}\quad C=W^{\mathrm{univ}}\quad\text{le plus \unsure{grd|gnd} poss.}
 \end{aligned}
\]
On devra donc prendre
\[
 TF=F\cap W_{\mathrm{univ}}=(W^{\mathrm{univ}})_*\cap W_{\mathrm{univ}},
\]
\deleted{\unsure{alors}} \inserted{et} cet ensemble est bien clos à droite \unsure{?}
\deleted{\unsure{pour} \illegible{4}} \inserted{par hypothèse.}
\deleted{\unsure{pour} \illegible{2}}

\uline{Corollaire} \inserted{Supposons les conditions b$'$) du th. Scholie \unsure{satisfaites}} Les conditions suivantes \inserted{a), b) sur $W$} sont
équivalentes (et impliquent \unsure{la|les} cond. du $3^\circ$ du th.-Scholie)
\deleted{a) On a \unsure{(et} aussi celles du $2^\circ$) du th. Scholie}
% archive sheet number: 29
